\section{Computations encoded by first-order sentences}\label{sec:computability}
Two reductions answer different questions. The first converts a successful run into a deduction, and an unsuccessful \emph{halting} run into a finite countermodel. The second converts an arbitrary accepting run into a finite model and recovers a run from every such model. Both reductions have arithmetical converses or separation consequences that are provable without deciding whether an arbitrary run halts.

\subsection{The numerical machine model}
Use deterministic machines with two stacks over a finite alphabet. An instruction either pushes a fixed letter on a selected stack and jumps, or inspects and pops a selected stack, with one jump for empty (leave it empty) and a designated jump for each popped letter. There are two stopping labels, yes and no. The input word starts on the right stack, its first letter uppermost, and the left stack starts empty. A program is a numbered finite table. A missing instruction can be treated as a permanent loop. The bounded-run predicate iterates the table lookup and the list updates for $t$ steps, freezing a stopped configuration. It is PR. Program codes and input words convert to each other by any fixed self-delimiting binary serialization.

Here are the coding facts needed for the diagonal arguments. A two-stack machine simulates a one-tape machine by keeping the cells left of the head, nearest first, on the left stack and the current cell and cells to its right on the right. An empty read returns blank. To move right, pop the old current cell and push its new value on the left. To move left, pop the old current cell, push its new value on the right, then transfer the nearest left cell (or a blank) onto the right. A stay is a right-stack top replacement. A finite instruction block performs each update and can remember any one letter and transition in its control. Conversely, ordinary finite scans, copying and the evaluation of any \emph{fixed} PR definition can be programmed on a tape. Store a list of binary registers separated by delimiters. A scan copies or replaces a register using a marked cursor, and successor is ordinary carry. To implement composition reserve disjoint scratch registers for the component computations. To implement primitive recursion copy the counter, compute the initial value and run the step block once for each counter decrement, passing the iteration index, previous value and untouched parameters. Induction on this fixed program construction gives PR running-time bounds for each PR routine: the recursion case iterates a PR upper bound for running time and register sizes together. Its iPRA invariant records the contents of the designated registers at block boundaries. A tape with auxiliary tracks uses a fixed product alphabet. These constructions justify precomposing a hypothetical machine with each particular compiler below, making a program that calls it on a transformed input, swaps its two stopping labels, or runs an unbounded enumeration one stage at a time. Only its finite traces are used in iPRA.

\subsection{A proof or a finite countermodel}
\begin{theorem}[Constructive Church separation]\label{thm:church}
There is a PR map $(P,w)\mapsto F_{P,w}$ to closed sentences. From a yes trace obtain an intuitionistic proof of $F_{P,w}$. From a no trace obtain a finite model of $\neg F_{P,w}$. No total computable bit test accepts every intuitionistically provable sentence and rejects every sentence with a finite Boolean countermodel. This holds already with equality and just one nonlogical binary relation.
\end{theorem}
\begin{proof}
First use a constant $e$, a unary constructor $c_a$ for each letter, a binary predicate $C_i$ for each program label, and a nullary predicate $Y$. The word term $[u]$ is $e$ for empty and $c_a([v])$ for $u=av$. Form a finite list $T$ of Horn sentences, including $C_{\rm start}(e,[w])$ and $\forall l,r\,(C_{\rm yes}(l,r)\to Y)$. For push-left from $i$ to $j$ add
\[
             C_i(l,r)\to C_j(c_a(l),r).
\]
For pop-left add $C_i(e,r)\to C_{j_0}(e,r)$ and
$C_i(c_a(l),r)\to C_{j_a}(l,r)$ for each letter $a$. Close all clauses universally. Right-stack rules exchange the two argument positions. Stop labels have no instruction clauses. Put $F=(\bigwedge T)\to Y$. Given an accepting trace, start with the initial fact, instantiate the clause for each instruction with the concrete word terms, and use modus ponens. The final clause yields $Y$. Conjunction projections and a single discharge produce the required proof.

For a rejecting trace let $B$ be the maximum length of any stack that occurs in it. Take all words of length at most $B$ and an overflow point $*$. Interpret $e$ as empty, $c_a$ as prefixing $a$ if the result fits, and as $*$ otherwise, including on $*$. Set $C_i$ to exactly the pairs of actual stack words visited at $i$. Make $Y$ false. Every push from a true antecedent is a step of the trace and fits the bound. The only preimage of a true nonempty word $av$ under $c_a$ is $v$. It has no preimage under another constructor. Thus every matching pop clause from a true antecedent produces the actual next pair. No antecedent at yes is true. These finite tables satisfy $T$ and falsify $F$, also when $B=0$.

Suppose a separator $D$ is given. Make a program $J$ that on input the code $u$ computes $F_{u,u}$ and runs $D$ on it, going to no on bit one and yes on bit zero. Take $u$ to be the code of $J$. Totality supplies the stopping trace. On bit zero the first construction supplies a proof that forces $D$ to output one. On bit one the second supplies a finite countermodel that forces zero. This contradiction is intuitionistic: the disjunction is read off a supplied terminating trace. Equivalently the arithmetic theorem says that no machine with a proof of total two-valued behavior can satisfy the two universally stated separator conditions, and, internally, negates the existence of a code satisfying their conjunction. The argument uses the test's claimed totality as a logical hypothesis.
\end{proof}

\paragraph{Reduction to a single binary relation.}
Here and below we use the following explicit signature encoding to justify the last sentence of the theorem. First replace each function, including a constant, by its graph relation under totality and uniqueness. Work on an inhabited sort $D$. A finite countermodel gives the actual graph tables. For the yes direction of the Horn reduction only a finite list of ground terms occurs in the constructed proof. Under the new graph axioms pick values successively for their subterms by totality. Use existential elimination to retain these fresh values locally and uniqueness to identify repeated occurrences. The translation of each original Horn axiom at those values gives precisely its ground consequence. More explicitly its antecedent universally quantifies over the named variable values. Its matrix says that for all consistent evaluations of the terms in that matrix the evaluated atoms satisfy the implication. For nullary facts use the same convention. The final atom $Y$ mentions none of the chosen values, so all the eliminations can be closed. This gives a PR intuitionistic proof compiler for the only direction used in Theorem~\ref{thm:church}.

For a finite relational signature use a directed edge $E$. Use parameter vertices $d$ (objects), $u_R$ (tuple nodes for relation $R$) and $p_{R,j}$ (coordinate nodes). A vertex $v$ has a stated type with tag $h$ if $E(v,h)$. Translate variables by relativizing quantifiers to $D(x):=E(x,d)$. Translate $R(x_1,\ldots,x_k)$ by
\[
 \exists v\left(E(v,u_R)\ \wedge \bigwedge_{1\le j\le k}\exists b\,
       (E(v,b)\wedge E(b,p_{R,j})\wedge E(b,x_j))\right).
\]
For $k=0$ the conjunction is empty. Equality is actual equality on $D$. Every graph and parameter tuple with inhabited $D$ now gives a decoded relational structure. Bounded scans give the truth identity by induction on formulas. For the converse encode any table by disjoint fresh vertices for all tags, all elements, every true tuple, and each of that tuple's coordinates. Draw just the specified type, incidence and value edges. Tag vertices have no outgoing edges. Each coordinate node links to its tag and its unique data element. Each tuple node links to its tag and its own coordinates. Among typed data the translation consequently recovers exactly the original tuples, including repetitions and nullaries.

For the separator use universal quantification over all tags, and an implication from inhabitance of $D$ and the relativized graph axioms and Horn assumptions to $Y^*$. A proof from an accepting trace runs verbatim inside the interpreted sort, using the inhabitance and totality hypotheses, then discharges and generalizes. In a rejecting graph the designated tags falsify that universal implication. For a finite-satisfiability reduction use existential tags, conjoin $\exists x\,D(x)$ and the translated sentence. The two decoding arguments give equivalence in both directions. Sorted inputs can first be relativized as in the finite-model chapter. All tags, formulas and finite incidence graphs are built by bounded loops.

\begin{corollary}
Intuitionistic and classical derivability are recursively enumerable and undecidable, even over the indicated binary signature. Finite validity is undecidable there as well. A decision procedure for ordinary classical validity, in its usual external reading, would likewise give the forbidden separator. There is no PR bound on the least proof code of a provable sentence as a function of its code.
\end{corollary}
\begin{proof}
The search in Section~\ref{sec:deduction} enumerates proofs. Finite soundness ensures that either proof calculus separates the yes and no cases if its theorem predicate were decidable. Finite validity separates them for the same reason. Ordinary semantics on all structures is sound for the displayed proofs and includes the finite countermodels. A PR upper bound, even one required only on provable inputs, would decide by bounded search for an accepted proof.
\end{proof}

\subsection{Recovering a real computation from a finite model}
\begin{theorem}[Finite satisfiability and Trakhtenbrot's obstruction]\label{thm:trakh}
One can compute a sentence $B_{P,w}$ in the language of a directed graph with equality with
\[
 \exists t\ \mathrm{Accept}(P,w,t)
 \quad\longleftrightarrow\quad
 \exists M\,[\mathrm{FinGraph}(M)\wedge V(M,B_{P,w})=1].
\]
There are PR witness maps both ways. Finite satisfiability is recursively enumerable, while finite validity has no recursive enumeration.
\end{theorem}
\begin{proof}
Describe the sentence first with two nonempty sorts, times and stack objects. Use a strict linear order on the times, a root $e$ of the stack sort, binary relations $U_a(v,v')$ for adding $a$ to stack $v$, two stack-valued functions $l(t),r(t)$, and a unary predicate $Q_i(t)$ for each label. Require the $Q_i$ to partition the times. Require $e$ to have no incoming $U$ edge and require every child to have at most one incoming parent-letter pair. Explicitly, $U_a(x,z)\wedge U_b(y,z)$ implies $x=y$ if $a=b$ and implies $\bot$ otherwise. The times $t_0,t_f$ are existentially specified least and greatest. The first label is start, $l(t_0)=e$, and $r(t_0)$ is reached from $e$ by a chain spelling the input in reverse order of its letters. Write this finite chain with existential intermediate objects. Require $Q_{\rm yes}(t_f)$.

Put $\mathrm{next}(t,t') := t<t'\wedge\neg\exists s(t<s\wedge s<t')$. For every adjacent pair require the disjunction of the legal instruction cases. For a push-left case its conjuncts are $Q_i(t),Q_j(t'),r(t')=r(t),U_a(l(t),l(t'))$. For pop-empty-left use the two corresponding labels, right-stack equality and $l(t)=e=l(t')$. For pop-$a$-left use its labels, right-stack equality and $U_a(l(t'),l(t))$. Include the symmetric right-stack cases. There are no disjuncts for transitions from a stopping label. These instructions expand to actual first-order formulas, of size bounded by a fixed polynomial in the written table and word if the two-sorted presentation is used.

An accepting run supplies time $0,\ldots,t$ up to its first stop, all its stack words and their suffixes, and the actual prefix edges between included words. Fill the two stack functions and each state bit from the trace. Every required conjunct follows from the successive configurations. In the converse direction sort the finite time table by its strict order: find the least unused element by bounded comparison, then mark it used. Every nonempty unused sublist has a least element, obtained by comparing its elements successively and retaining the smaller. Its successive choices exhaust the times and each consecutive pair satisfies $\mathrm{next}$. Carry, for each of the two current stack words, the list of object vertices from $e$ to the indicated top together with the letters. Initialization gives the first lists. A push appends the displayed edge. A pop from an empty word can only use the empty case because $e$ has no incoming edge. For a nonempty word the top has a recorded incoming edge, so differs from $e$, and uniqueness of the incoming parent and letter returns precisely the recorded predecessor. Delete that edge. Thus each move is precisely a machine step, and the last label certifies acceptance. Neither cycles outside the used paths nor unnamed stack objects affect the extraction. The induction stores words of length at most $|w|+$ the number of times and the two path lists. This is a bounded bit invariant of the time scan. Translate the functions, sorts and finite relational tables to a single graph by the preceding encoding. Reconstruction from graph to the time trace is a composition of PR operations.

Enumerating finite graph codes and evaluating a sentence semidecides finite satisfiability. For the final assertion suppose a machine enumerates exactly the finitely valid graph sentences, with an eventual output for each such sentence. Build $L$ that on input $u$ scans the enumeration for the sentence $\neg B_{u,u}$ and halts at yes just when it finds it. Set $u$ equal to the code of $L$ and write $H$ for acceptance of this run. By the construction and the assumed enumeration,
\[
 H\ \longleftrightarrow\ \neg B_{u,u}\text{ is finitely valid}
       \ \longleftrightarrow\ \neg H .
\]
The second equivalence uses finite Boolean evaluation and the two witness maps. From $H\to\neg H$ one gets $\neg H$, and then the other direction gives $H$, a contradiction. No case split on an unresolved halting question occurs.
\end{proof}
